\section{Deep Learning e Redes Neurais}

\textbf{OBJETIVO GERAL:}

Desenvolver competências para projetar, treinar, avaliar e implantar modelos de Deep Learning, utilizando redes neurais artificiais e arquiteturas modernas para resolução de problemas complexos envolvendo classificação, regressão, visão computacional, processamento de linguagem natural e Inteligência Artificial Generativa.

\textbf{CARGA HORÁRIA:} 160 horas

\textbf{FUNÇÃO:}

Desenvolver sistemas computacionais baseados em Inteligência Artificial, Análise Preditiva e Ecossistemas IoT, atendendo às normas de qualidade corporativas, robustez metodológica, governança de dados e arquitetura segura.

\textbf{SUBFUNÇÕES:}

\begin{itemize}
\item Preparar bases de dados para treinamento de redes neurais.
\item Desenvolver modelos de Deep Learning.
\item Avaliar desempenho de redes neurais profundas.
\item Otimizar modelos utilizando técnicas modernas de treinamento.
\item Implantar modelos em aplicações computacionais.
\end{itemize}

\textbf{PADRÕES DE DESEMPENHO:}

\begin{itemize}
\item Preparar conjuntos de dados para treinamento.
\item Configurar arquiteturas de redes neurais.
\item Treinar modelos utilizando GPUs.
\item Avaliar desempenho utilizando métricas apropriadas.
\item Ajustar hiperparâmetros.
\item Aplicar técnicas para redução de overfitting.
\item Utilizar aprendizado por transferência.
\item Implantar modelos em aplicações práticas.
\end{itemize}

\textbf{CAPACIDADES TÉCNICAS:}

\textbf{Fundamentos de Redes Neurais}

\begin{itemize}
\item Interpretar a arquitetura de redes neurais artificiais.
\item Compreender o funcionamento dos neurônios artificiais.
\item Diferenciar arquiteturas de Deep Learning.
\item Relacionar funções de ativação aos modelos.
\end{itemize}

\textbf{Treinamento de Modelos}

\begin{itemize}
\item Preparar bases de treinamento.
\item Configurar hiperparâmetros.
\item Executar processos de treinamento.
\item Validar modelos durante o treinamento.
\end{itemize}

\textbf{Otimização}

\begin{itemize}
\item Aplicar algoritmos de otimização.
\item Utilizar técnicas de regularização.
\item Reduzir overfitting.
\item Melhorar desempenho dos modelos.
\end{itemize}

\textbf{Transfer Learning}

\begin{itemize}
\item Utilizar modelos pré-treinados.
\item Adaptar modelos para novos problemas.
\item Realizar fine-tuning.
\item Avaliar desempenho após adaptação.
\end{itemize}

\textbf{Implantação}

\begin{itemize}
\item Exportar modelos treinados.
\item Integrar modelos em aplicações.
\item Avaliar desempenho em produção.
\item Documentar modelos desenvolvidos.
\end{itemize}

\textbf{CONHECIMENTOS:}

\textbf{Fundamentos de Deep Learning}

\begin{itemize}
\item Redes Neurais Artificiais
\item Neurônios
\item Camadas
\item Pesos
\item Bias
\item Forward Propagation
\item Backpropagation
\end{itemize}

\textbf{Funções de Ativação}

\begin{itemize}
\item Sigmoid
\item Tanh
\item ReLU
\item Leaky ReLU
\item Softmax
\end{itemize}

\textbf{Treinamento}

\begin{itemize}
\item Epochs
\item Batch Size
\item Learning Rate
\item Funções de Perda
\item Gradiente Descendente
\end{itemize}

\textbf{Otimizadores}

\begin{itemize}
\item SGD
\item Momentum
\item RMSProp
\item Adam
\item AdamW
\end{itemize}

\textbf{Regularização}

\begin{itemize}
\item Dropout
\item Batch Normalization
\item Early Stopping
\item L1
\item L2
\end{itemize}

\textbf{Arquiteturas}

\begin{itemize}
\item Perceptron Multicamadas
\item Redes Convolucionais (CNN)
\item Redes Recorrentes (RNN)
\item LSTM
\item GRU
\item Transformers (introdução)
\end{itemize}

\textbf{Transfer Learning}

\begin{itemize}
\item Modelos Pré-treinados
\item Fine-Tuning
\item Feature Extraction
\item Congelamento de Camadas
\end{itemize}

\textbf{Frameworks}

\begin{itemize}
\item TensorFlow
\item Keras
\item PyTorch
\item Hugging Face (introdução)
\end{itemize}

\textbf{Avaliação}

\begin{itemize}
\item Curvas de Aprendizado
\item Overfitting
\item Underfitting
\item Métricas de Desempenho
\item Ajuste de Hiperparâmetros
\end{itemize}

\textbf{Implantação}

\begin{itemize}
\item Exportação de Modelos
\item Inferência
\item APIs para Modelos
\item Deploy em Nuvem
\end{itemize}

\textbf{CAPACIDADES SOCIOEMOCIONAIS:}

\begin{itemize}
\item Demonstrar persistência durante o treinamento de modelos complexos.
\item Avaliar criticamente os resultados obtidos.
\item Formular hipóteses para melhoria contínua dos modelos.
\item Trabalhar colaborativamente em projetos de Inteligência Artificial.
\item Produzir documentação técnica dos experimentos realizados.
\item Compartilhar conhecimento sobre arquiteturas de Deep Learning.
\item Demonstrar organização na condução dos experimentos.
\item Agir com responsabilidade na utilização de modelos inteligentes.
\end{itemize}

\textbf{AMBIENTES PEDAGÓGICOS:}

\begin{itemize}
\item Laboratório de Inteligência Artificial.
\item Laboratório de Ciência de Dados.
\item Ambiente Virtual de Aprendizagem.
\item Ambiente Cloud com suporte a GPU.
\end{itemize}

\textbf{EQUIPAMENTOS E RECURSOS:}

\textbf{Hardware}

\begin{itemize}
\item Computadores com acesso à Internet.
\item Estações com GPU para treinamento de modelos.
\item Projetor multimídia.
\end{itemize}

\textbf{Software}

\begin{itemize}
\item Python
\item TensorFlow
\item Keras
\item PyTorch
\item Jupyter Notebook
\item Google Colab
\item VS Code
\item Git
\item Docker
\end{itemize}
